\setauthor{Markus Tran}
Die vorliegende Anwendung adressiert ein konkretes Problem des modernen
Medienkonsums: die zunehmende Unübersichtlichkeit bei der Verwaltung mehrerer
paralleler Streaming-Abonnements. Um die Relevanz und den Mehrwert der
entwickelten Lösung zu verdeutlichen, werden in diesem Kapitel die fachlichen
Grundlagen des Abo- und Mediatheken-Managements erläutert.


\section{Zentrale Verwaltung von Abonnements und Laufzeiten}

Der Markt für Video-on-Demand-Dienste hat sich in den vergangenen Jahren stark
ausdifferenziert. Neben den frühen Marktführern wie Netflix und Amazon Prime
Video sind zahlreiche weitere Anbieter entstanden -- darunter Disney+, Apple TV+,
Paramount+ und viele nationale Mediatheken. Eine repräsentative Studie des
Digitalverbands Bitkom aus dem Jahr 2023 zeigt, dass ein Großteil der deutschen
Haushalte gleichzeitig mehrere Streaming-Dienste abonniert hat.

Diese Fragmentierung des Angebots führt zu einer praktischen Herausforderung:
Nutzerinnen und Nutzer verlieren den Überblick über ihre aktiven Abonnements,
die jeweiligen Abrechnungszeitpunkte und die damit verbundenen monatlichen oder
jährlichen Kosten. Anders als bei klassischen Dauerschuldverhältnissen -- etwa
Mobilfunkverträgen -- erfolgt die Abrechnung von Streaming-Diensten häufig
automatisch und ohne gesonderte Zahlungsbenachrichtigung, was eine unbewusste
Fortsetzung nicht mehr genutzter Abonnements begünstigt.

Eine zentrale Verwaltung von Abonnements und deren Laufzeiten schafft hier
Abhilfe. Durch die strukturierte Erfassung aller aktiven Dienste, ihrer
Abrechnungszyklen -- monatlich oder jährlich -- sowie der jeweiligen
Abrechnungsdaten erhalten Nutzer einen konsolidierten Überblick über ihre
laufenden Verpflichtungen. Dies ermöglicht eine bewusste Entscheidung darüber,
welche Dienste tatsächlich genutzt werden und welche gekündigt werden könnten.


\section{Darstellung von Verfügbarkeiten von Filmen und Serien}

Ein weiteres zentrales Problem des fragmentierten Streaming-Markts ist die
uneinheitliche Verfügbarkeit von Inhalten. Ein Film oder eine Serie ist
selten auf allen Plattformen gleichzeitig verfügbar. Stattdessen sind Inhalte
exklusiv an bestimmte Anbieter gebunden oder wechseln im Laufe der Zeit
zwischen verschiedenen Plattformen.

Für Nutzer bedeutet dies, dass die Suche nach einem konkreten Titel häufig
mehrere manuelle Suchanfragen auf unterschiedlichen Plattformen erfordert.
Dieser Prozess ist zeitaufwändig und frustrierend, insbesondere wenn der
gesuchte Inhalt auf keiner der abonnierten Plattformen verfügbar ist.

Eine Anwendung, die Filminformationen mit der Verfügbarkeit bei
Streaming-Anbietern verknüpft, löst dieses Problem. Nutzer können direkt
bei der Filmsuche erkennen, auf welchen Plattformen ein Titel verfügbar
ist und -- in Kombination mit der Abonnementverwaltung -- ob dieser Titel
auf einer ihrer bereits abonnierten Plattformen abrufbar ist. Diese
Verknüpfung von Inhaltsdaten und Abonnementstatus stellt den zentralen
fachlichen Mehrwert der vorliegenden Anwendung dar.

Die technische Grundlage für diese Funktionalität bildet die TMDB API,
die neben umfangreichen Filmdaten auch länderspezifische Informationen
zu Streaming-Verfügbarkeiten bereitstellt. Dabei wird zwischen
verschiedenen Bezugsmodellen unterschieden: Flatrate-Verfügbarkeit
im Rahmen eines Abonnements, Leihmodelle sowie Kaufoptionen. Im Kontext
der vorliegenden Anwendung sind ausschließlich Flatrate-Angebote relevant,
da der Fokus auf der Verwaltung aktiver Abonnements liegt.


\section{Mehrwert einer zentralen Applikation für Nutzer}

Der Mehrwert der entwickelten Anwendung ergibt sich aus der Zusammenführung
zweier bisher getrennter Informationsquellen: der persönlichen
Abonnementverwaltung und der Suche nach Streaming-Verfügbarkeiten. Bestehende
Lösungen adressieren diese Bereiche typischerweise isoliert voneinander.

Allgemeine Budgetverwaltungs-Apps ermöglichen die Erfassung von
Abonnementkosten, bieten jedoch keine Verknüpfung mit konkreten Filminhalten.
Streaming-Suchmaschinen wie JustWatch zeigen an, auf welchen Plattformen ein
Film verfügbar ist, verfügen jedoch über keine persönliche
Abonnementverwaltung, die den individuellen Nutzerstatus berücksichtigt.

Die vorliegende Anwendung kombiniert beide Aspekte in einer einheitlichen
Oberfläche. Nutzer können ihre Abonnements verwalten, nach Filmen suchen
und direkt erkennen, ob ein gesuchter Titel auf einer ihrer abonnierten
Plattformen verfügbar ist. Dieser personalisierte Ansatz reduziert den
kognitiven Aufwand bei Entscheidungen über Streaming-Inhalte und schafft
einen konkreten praktischen Nutzen im Alltag.

Darüber hinaus trägt die Anwendung zur Kostentransparenz bei. Durch die
strukturierte Erfassung von Abrechnungszyklen und -daten können Nutzer
ihre monatlichen Ausgaben für Streaming-Dienste besser nachvollziehen
und bewusster gestalten. Gerade bei jährlichen Abrechnungsmodellen, die
aufgrund ihres günstigeren Preises häufig gewählt werden, geht das
konkrete Abrechnungsdatum im Alltag leicht verloren.


\section{Datenschutz- und Sicherheitsanforderungen}

Da die vorliegende Anwendung personenbezogene Daten verarbeitet --
insbesondere Informationen über abonnierte Streaming-Dienste und deren
Abrechnungsdaten -- sind Datenschutz- und Sicherheitsanforderungen ein
wesentlicher Bestandteil des Systemdesigns.

Im Sinne der Datensparsamkeit, einem Grundprinzip der Datenschutz-Grundverordnung
(DSGVO), werden ausschließlich jene Daten erfasst und gespeichert, die für
die Kernfunktionalität der Anwendung erforderlich sind. Auf die Speicherung
besonders sensibler Daten wie Zahlungsinformationen oder Zugangsdaten zu
den jeweiligen Streaming-Plattformen wird bewusst verzichtet. Die Anwendung
beschränkt sich auf die Verwaltung von Metadaten der Abonnements --
Anbieter, Abrechnungszyklus und Datum der letzten Abbuchung.

Die Authentifizierung und Autorisierung der Nutzer wird vollständig über
Keycloak abgewickelt, ein etabliertes Open-Source-System für
Identity- und Access-Management. Durch diese Auslagerung wird das Risiko
einer unsicheren Eigenimplementierung von Passwortlogik eliminiert.
Zugangsdaten werden zu keinem Zeitpunkt durch die Anwendung selbst
verarbeitet oder gespeichert.

Der Zugriff auf benutzerspezifische Daten ist durch eine
rollenbasierte Zugriffskontrolle abgesichert. Geschützte Endpunkte des
Backends sind ausschließlich für authentifizierte Nutzer mit den
entsprechenden Rollen zugänglich. Das Frontend steuert darüber hinaus
den Zugriff auf geschützte Navigationsbereiche, sodass persönliche
Inhalte nicht ohne gültige Anmeldung angezeigt werden. Die Kommunikation
zwischen Frontend, Backend und Keycloak erfolgt ausschließlich über
verschlüsselte HTTPS-Verbindungen.